\section{课程定位与核心问题}

\subsection{为什么这讲重要}
这节课并不是再介绍一个 ``更大模型''，而是把注意力放到一个更工程化的问题：当模型进入真实网页环境，如何从 ``会回答问题'' 变成 ``会完成任务''。Ruslan 把主题压缩成三件事：评估、搜索、数据扩展。这三件事恰好是大多数 Agent 系统落地时最先撞墙的地方。

\begin{importantbox}{本讲主线}
\begin{enumerate}
    \item \textbf{VisualWebArena}：如何在可复现环境里评估多模态 Web Agent；
    \item \textbf{Inference-Time Search}：如何在推理时让 Agent 更会试错和回溯；
    \item \textbf{Data Scaling}：如何在数据稀缺条件下持续提升 Agent 能力。
\end{enumerate}
\end{importantbox}

\begin{knowledgebox}{从 ``聊天'' 到 ``执行'' 的范式变化}
当任务从问答变成网页操作时，模型要同时处理视觉输入、状态变化、动作约束、长期目标和失败恢复。这个问题更接近 ``Sequential Decision Making''，而不是单轮文本生成。
\end{knowledgebox}

\begin{figure}[H]
\centering
\includegraphics[width=0.82\textwidth]{frame-01-visualwebarena.jpg}
\caption{视频约 00:06:50：讲者引入 VisualWebArena，强调在真实交互环境评估多模态 Agent}
\end{figure}

\subsection{这门课对实践者的直接启发}
讲者反复强调一点：Agent 系统的瓶颈不只在模型参数，而在系统化能力。即便基础模型更强，如果没有搜索、验证器和高质量交互数据，性能上限仍然会很快触顶。反过来，哪怕基础模型不是最强，只要在这三条链路上做对，也可能拿到可观收益。

\begin{warningbox}{常见误区}
把 ``模型能力提升'' 直接等同于 ``任务完成率提升'' 往往是错的。Web 任务里，状态空间大、步骤长、中间错误代价高，单步能力提升并不会自动转化为长链路成功率。
\end{warningbox}

\subsection{本章小结}
本讲聚焦 Agent 工程的三根支柱：评估、搜索、数据。它关注的不是单次回答质量，而是长链路任务的成功率与可扩展性。

